\subsection{提升定理}

命题 1．——设 \(k_{0}\) 是一个域，A 是一个 \(k_{0}\)-代数且是分离完备的局部环，K 是一个 \(k_{0}\)-子扩张，属于 \(\kappa_{\mathrm{A}}\)，并且具有可分超越基 \((\xi_{\lambda})_{\lambda \in \Lambda}\)，它在 \(k_{0}\) 上（A，V，p. 130，定义 1）。对于每个 \(\lambda \in \Lambda\)，令 \(x_{\lambda}\) 为 A 中 \(\xi_{\lambda}\) 的一个代表元。存在 A 的唯一子域 L，包含 \(k_{0}\) 和元素 \(x_{\lambda}\)，使典范同态 \(\pi\) 从 A 到 \(\kappa_{\mathrm{A}}\) 诱导出从 L 到 K 的同构。

设 \(\varphi\) 为一个 \(k_{0}\)-同态，从多项式环 \(k_{0}[(X_{\lambda})_{\lambda \in \Lambda}]\) 到 A，将 \(X_{\lambda}\) 映到 \(x_{\lambda}\)，对每个 \(\lambda \in \Lambda\) 均如此。设 \(u\) 是 \(k_{0}[(X_{\lambda})_{\lambda \in \Lambda}]\) 的非零元素；有 \(\pi (\varphi (u)) \neq 0\)，因为族 \((\xi_{\lambda})_{\lambda \in \Lambda}\) 关于 \(k_{0}\) 在 \(\kappa_{\mathbf{A}}\) 中代数自由；因而 \(\varphi (u)\) 在局部环 A 中可逆。因此 \(\varphi\) 延拓为同态 \(\psi\)，从域 \(k_{1} = k_{0}((X_{\lambda})_{\lambda \in \Lambda})\) 到 A。于是 A 是一个 \(k_{1}\)-代数，\(\kappa_{\mathbf{A}}\) 是 \(k_{1}\) 的扩张，而 K 是 \(\kappa_{\mathbf{A}}\) 的一个子扩张，它关于 \(k_{1}\) 代数且可分。还需证明存在 A 的唯一子域 L，包含 \(\psi(k_{1})\)，并满足 \(\pi(L) = K\)。

\begin{enumerate}
\def\labelenumi{\alph{enumi})}
\item
  存在性：令 S 为 A 中包含 \(\psi(k_1)\) 且满足 \(\pi(L) \subset K\) 的子域所成的集合；它关于包含关系是归纳的。令 L 为 S 的极大元素；把 K 看作 L 的代数可分扩张（A，V，p. 40，命题 9）。令 \(\xi \in K\)，并令 \(P \in L[X]\) 为其在 L 上的最小多项式。由于 \(\xi\) 是 P 的单根，Hensel 引理（III，§ 4，第 5 号，定理 2 的推论 1）保证存在 A 中的元素 \(x\)，使 \(\pi(x) = \xi\) 且 \(P(x) = 0\)。A 的子环 \(L[X]\) 属于 S；根据 L 的极大性，有 \(x \in L\)，从而 \(\xi \in \pi(L)\)。最后，有 \(\pi(L) = K\)。
\item
  唯一性：设 L 和 L' 是 A 中包含 \(\psi(k_1)\) 且满足 \(\pi(L) = \pi(L') = K\) 的两个子域。令 \(\xi \in K\)，并令 \(x \in L\) 和 \(x' \in L'\) 为满足 \(\pi(x) = \pi(x') = \xi\) 的元素。若 \(P \in k_1[X]\) 是 \(\xi\) 在 \(k_1\) 上的最小多项式，则 \(\xi\) 是 P 的单根，且有 \(P(x) = P(x') = 0\)。根据 Hensel 引理（同上），有 \(x = x'\)。所以有 \(L = L'\)。
\end{enumerate}

注记．——* 上面的证明更一般地适用于只假设 A 是 Hensel 局部环的情形*。唯一性证明使用了局部环 A 分离这一假设，但没有使用其完备性。

